\chapter*{Introduction Generale}
\phantomsection
\addcontentsline{toc}{chapter}{Introduction Generale}


Dans un environnement industriel de plus en plus exigeant, la maîtrise des flux logistiques représente un facteur essentiel pour assurer la continuité de la production, optimiser les opérations et améliorer la prise de décision. Les entreprises industrielles sont amenées à gérer plusieurs mouvements de transport, qu'il s'agisse de l'approvisionnement en matières premières ou de la livraison de produits finis. Dans ce contexte, la digitalisation des processus permet de renforcer la traçabilité, de réduire la dépendance aux échanges manuels et de centraliser les informations nécessaires au suivi des opérations.

Le présent projet s'inscrit dans le cadre d'un stage de fin d'études réalisé au sein de SOMASTEEL, une entreprise marocaine opérant dans le secteur de la sidérurgie. L'entreprise est spécialisée dans le laminage à chaud des aciers et la production de fer à béton. Son activité industrielle repose notamment sur l'approvisionnement en ferraille, matière première utilisée dans le processus de production, ainsi que sur la gestion des flux de transport associés.

SOMASTEEL dispose de deux catégories de camions. La première catégorie regroupe les camions appartenant à l'entreprise, utilisés principalement pour la livraison des produits finis. Ces camions sont équipés d'un système GPS permettant de suivre leurs déplacements. La deuxième catégorie concerne les camions loués, utilisés pour transporter la ferraille importée depuis le port vers le site industriel de l'entreprise. Cependant, contrairement aux camions internes, ces véhicules loués ne peuvent pas être équipés de dispositifs GPS, car l'entreprise ne dispose pas de l'autorisation nécessaire pour installer ce type de matériel sur des camions qui ne lui appartiennent pas.

Cette contrainte a mis en évidence un besoin important au niveau du service logistique : disposer d'une solution alternative permettant de suivre les camions loués entre le port et l'entreprise. Avant la mise en place du projet, le suivi de ces camions reposait principalement sur des échanges téléphoniques, ce qui rendait l'information moins centralisée, moins précise et plus difficile à exploiter. Le projet vise donc à proposer une solution digitale capable d'assurer le suivi des expéditions sans recourir au GPS.

La solution développée, intitulée \textbf{SomaScan}, repose sur l'utilisation de codes QR associés aux plaques d'immatriculation des camions. Chaque camion est identifié par un code QR unique, généré automatiquement lors de son ajout dans le système. Deux opérateurs interviennent sur le terrain : un opérateur au niveau du port et un opérateur au niveau de l'usine. Chacun scanne le code QR du camion afin d'enregistrer son passage et de mettre à jour son statut dans le système. Un agent logistique dispose, quant à lui, d'une interface web d'administration lui permettant de gérer les camions, suivre les expéditions et préparer des rapports destinés à la Direction Générale.

Sur le plan technique, la solution est composée de trois applications principales. Le backend est développé avec Laravel et expose des API REST assurant la gestion des données, de l'authentification, des rôles et de la logique métier. L'application mobile est développée avec Flutter et permet aux opérateurs de scanner les codes QR sur le terrain. L'application web est développée avec React et fournit à l'agent logistique une interface d'administration et de suivi.

Ce rapport présente l'ensemble du travail réalisé durant le stage, depuis l'analyse du besoin jusqu'à la conception, le développement, les tests, le déploiement et la proposition d'un workflow DevOps. Il est structuré en cinq chapitres. Le premier chapitre présente le contexte général du projet, l'organisme d'accueil, la problématique et les objectifs. Le deuxième chapitre est consacré à l'analyse et à la spécification des besoins. Le troisième chapitre décrit la conception et l'architecture de la solution. Le quatrième chapitre présente la phase d'implémentation. Enfin, le cinquième chapitre aborde le déploiement de l'application ainsi que l'approche DevOps proposée pour automatiser et fiabiliser les futures mises en production.
